% 性能ベンチマーク：和文主体の長文（禁則・約物・混植を多数含む）
% 「通常の短い記事（和文主体）」と「約物・短段落・複雑な行分割を
% 多数含むストレス文書」を兼ねる入力。
\documentclass[pdflatex]{jlreq}
\begin{document}
\section{序論}
日本語組版の規則は，文字クラスとその並びによって空き量を定める。例えば，
「括弧の前後」には二分アキを入れる（3.1.5節）。また、句読点が行頭に来る
ことは禁則により許されない。English text is mixed into paragraphs like this
one, with TeX typesetting both scripts. 縦組みの話はここでは扱わない。

「開き括弧に続く文章。閉じ括弧」の後、中点「・」や感嘆符「！」・「？」
のような区切り約物（クラス4）も現れる。脚注\footnote{これは脚注である。
規則（4.2.5）に従う。}や，数式 $a^2+b^2=c^2$ も混在する。

あいうえおかきくけこさしすせそたちつてとなにぬねのはひふへほまみむめもや
ゆよらりるれろわをん。息の長い一文がこのように続く場合，行分割の候補数は
増加し，禁則ペナルティの評価回数も増える。001，002，…，099のような半角も
混ざる。これを数千字規模で繰り返すことで，文字処理経路の実行時間を測る。

\section{本論}
\subsection{規則の適用}
行頭禁則・行末禁則は，改行位置候補の直前・直後の文字クラスで決まる。
ペナルティ10000は「この位置での改行を禁止する」ことを意味する。
二分アキ・四分アキはglueとして挿入され，行の追込み（詰め）に使われる。

\subsection{測定条件}
ウォームアップ後の実行時間，生成PDFサイズ，ログの警告件数を比較する。
比較対象は(A) pdfLaTeX＋本バックエロード，(B) 同＋通常動作，
(C) LuaLaTeX＋既存jlreq，(D) upLaTeX＋既存jlreq。

% ---- 長文の繰り返し（機械生成ではなく固定テキストで決定的に） ----
\def\para#1{#1 #1 #1 #1 #1}
\def\parablock{日本語組版の規則は，文字クラスとその並びによって空き量を
定める。句読点「，」「。」が行頭に来ることは許されない。「括弧」・中点
「・」・引用'英語'も混在する。禁則，追込み，追出しの順で評価される。}

\parablock\parablock\parablock\parablock\parablock\parablock\parablock\parablock

\section{結論}
本ベンチマークは，約物・短段落・混植を多数含むストレス入力として機能する。
\end{document}
